% Diese Vorlage wurde von Benedikt Loepp
% im März 2018 aktualisiert.
% Version: 1.1

\documentclass[ngerman, a4paper, abstract=false, twoside=false, parskip=half, DIV=12, bibtotoc]{scrartcl}

\usepackage[T1]{fontenc}
\usepackage[utf8]{inputenc}

% Hier den eigenen Namen, die Matrikelnummer und den (vorläufigen)
% Titel der Abschlussarbeit eintragen
\newcommand{\autor}{Claudia Schwenger}
\newcommand{\matrikelnr}{2285577}
\newcommand{\titel}{Bursting the Bubble: Stärkung von Übersicht, Transparenz und Kontrolle in Empfehlungssystemen für Online-Nachrichten}

% Weitere Pakete können problemlos hinzugefügt werden.

% Neue deutsche Rechtschreibung
\usepackage[ngerman]{babel}

\usepackage{csquotes}

% Umbrüche in Zitaten erlauben
\usepackage{cite}

% Zitierungen a la 
% Ziegler (2010) durch citet{ziegler:2010}
% --> t wie textual
% bzw. (Ziegler, 2010) durch \citep{ziegler:2010} 
% --> p wie paranthesis
\usepackage{apacite}
\usepackage{natbib}

% Einbinden von Grafiken
\usepackage{graphicx}

% Komfortableres Tabellen-Paket
\usepackage{tabularx}

% Schriftfamilie Times
\usepackage{palatino}

%Schriftfarbe
\usepackage{xcolor}

% Ermöglicht mehrere Bilder in einer floatenden Figure-Umgebung
\usepackage{subfig}

% Erweiterte Mathe-Zeichen
\usepackage{amsmath,amsmath,amssymb,amsfonts}

% Erlauben von urls
\usepackage{url}

% Einbinden von Listings
\usepackage{listings}

% Klickbare Links innerhalb des PDFs
\usepackage[
	bookmarks=true,         % show bookmarks bar?
    unicode=false,          % non-Latin characters in Acrobat’s bookmarks
    pdftoolbar=true,        % show Acrobat’s toolbar?
    pdfmenubar=true,        % show Acrobat’s menu?
    pdffitwindow=false,     % window fit to page when opened
    pdfstartview={Fit},    % fits the width of the page to the window
    pdftitle={Proposal},    % title
    pdfauthor={\autor},     % author
    pdfsubject={\titel},   % subject of the document
    hidelinks=true,			  % do not show colored links
    pdfdisplaydoctitle=true,	  % show document title instead of filename in titlebar
    pdfpagelayout=TwoColumnRight	% show two pages by default
]{hyperref}

% Anpassung des Abstracts
\usepackage{xpatch}
\xpretocmd{\abstract}{\begin{addmargin}{1.5cm}}{}{}
\xpretocmd{\endabstract}{\end{addmargin}}{}{}
\xapptocmd{\abstract}{\noindent}{}{}
\setcounter{tocdepth}{2}
\setcounter{secnumdepth}{1}

\title{Proposal}
\subtitle{\titel}
\author{\autor \\ {\small Matrikelnr.: \matrikelnr}}
\date{\today}

\begin{document}
\maketitle

\begin{abstract}% Hier den Abstract einfügen
Kurze Zusammenfassung der geplanten Arbeit in 5-10 Zeilen. Ein Leser, der nur diese Zusammenfassung liest, sollte in der Lage sein zu verstehen was in der Arbeit geplant ist. Auch sollte kurz darauf eingegangen werden, wie sich die Thematik aus dem Stand der Forschung motiviert.
\end{abstract}

% Hier den Inhalt einfügen

\section{Problembeschreibung und Forschungsfrage}
\textcolor{red}{Kurzer Einführungstext, der das allgemeine Gebiet der Arbeit und dessen Hintergründe bündig beschreibt. Unter Verwendung geeigneter Literaturverweise sollte dabei klar werden, wo die Lücke in der bisherigen Forschung zu dem Thema liegt}.

-----------------------------------------------

Mit dem Einzug des digitalen Zeitalters sind die Möglichkeiten Informationen zu beschaffen immens gestiegen. Doch die verfügbare Menge an Informationen bringt auch ihre Schattenseite mit sich: Die immer weiter anwachsende Fülle an Informationen und Alternativen überfordert die Menschen in ihren Entscheidungsprozessen (\textcolor{red}{Richthammer & Pernul, 2016}) und der zugrundeliegende Information Overload (\textcolor{red}{evtl ergänzen})  gilt als typisches Problem unserer Zeit (\textcolor{red}{Jameson, 2003}). Um Menschen dabei zu helfen die benötigte Information zu finden und diese so darzustellen damit sie damit umgehen können (\textcolor{red}{Jameson, 2003}), wurden in der Vergangenheit zunehmend gezielt adaptive Systeme eingesetzt. Die automatische Filterung und Selektion soll dem Benutzer die Suche nach relevanten Informationen erleichtern, indem sie die Informationsfülle nach verschiedenen Kriterien vorfiltern und daraufhin bestimmte Objekte vorschlagen, die für den Benutzer voraussichtlich relevant sind \textcolor{red}{Richthammer & Pernul, 2016}). Diese Art von Systemen wird in der Regel als Empfehlungssystem bezeichnet, da sie auf Basis von Informationen über den Benutzer personalisiert Inhalte vorschlagen (\textcolor{red}{Haim, Graefe \&Brosius, 2017}). Die Informationen über den Benutzer, beispielsweise seine Interessen, seine Präferenzen oder sein Surfverhalten, werden explizit oder implizit (\textcolor{red}{Zuiderveen Borgesius et al., 2016}) gesammelt. Um Empfehlungen auf Grundlage von expliziten Informationen über den Benutzer zu generieren, muss dieser selbst aktiv werden und dem System seine Präferenzen mitteilen (zum Beispiel durch die Bewertung eines Produktes). Wird das, durch Beobachtung abgeleitete, Nutzungsverhalten als Grundlage für die Empfehlungen genutzt, handelt es sich um implizite Informationen. Häufig werden beide Informationsarten zur Empfehlungsgenerierung genutzt (\textcolor{red}{Haim, Graefe \& Brosius, 2017}). Dem Benutzer bleiben die Informationen, die dem System über ihn vorliegen und das daraus erstellte Nutzerprofil jedoch häufig verborgen \textcolor{red}{(Bakalov 2013)}. Daher kann oft nicht nachvollzogen werden, auf welcher Grundlage ein Empfehlungssystem die angezeigten Vorschläge generiert hat. In diesem Zusammenhang werden Empfehlungssysteme häufig als „Black-Box“ \textcolor{red}{(Quelle)} bezeichnet: Die Benutzer können die personalisierten Inhalte zwar sehen, jedoch keinen direkten Zugriff auf die Daten und keine Kontrolle über den Personalisierungsvorgang ausüben \textcolor{red}{(Bakalov 2013)}. Dies kann Misstrauen und Einbußen im Vertrauen hervorrufen \textcolor{red}{(Quelle)}. Es ist darüber hinaus sinnvoll den Nutzern Kontrolle über die Personalisierung zu erlauben, da sich dies positiv auf die wahrgenommene Nützlichkeit des Systems, die Bedienbarkeit und die Zufriedenheit der Nutzer mit dem personalisierten System auswirkt \textcolor{red}{(Bakalov 2013)}. 
Ein weiteres Problem, welches sich im Zusammenhang mit Empfehlungssystemen ergibt, ist, dass durch die automatische oder manuell veranlasste Filterung der Fokus häufig auf einen kleinen Ausschnitt des Informationsraumes gelenkt wird. Dadurch kann es passieren, dass verschiedene Alternativen herausfallen, die möglicherweise bedenkenswert wären, aber außerhalb der Filterkriterien liegen \textcolor{red}{(Richthammer & Pernul, 2016)}. 

Dies gilt auch für Online-Nachrichten: Der digitale Wandel ist auch im Bereich des Nachrichtenkonsums sichtbar: Es werden zunehmend Nachrichten aus Online Quellen genutzt. Eine umfassende Studie des Reuter Instituts mit 53.000 befragten Personen weltweit ergab, dass im Jahr 2016 bereits 23\% der Befragten Nachrichten hauptsächlich aus Online Quellen bezogen und weitere 44\% traditionellen (TV, Zeitung, usw.) und Online Medien die gleiche Bedeutung beimessen \textcolor{red}{(Newman, Fletcher, Levy, & Nielsen, 2016)}. Auch in dieser Domäne spielen Algorithmen eine wichtige Rolle. Ein Großteil der Online Konsumenten von Nachrichten gelangt über Suchmaschinen zu den gelesenen Artikeln (ca. 40\%) gefolgt von sozialen Medien (ca. 34\%) und Nachrichten-Aggregatoren (ca. 12\%) \textcolor{red}{(Newman, Fletcher, Levy, & Nielsen, 2016, s. 93)}. Auch hier wird angenommen, dass Empfehlungssysteme die Diversität der Nachrichten durch die Personalisierung und die damit verbundene individuelle Filterung entsprechend der Interessen des Benutzers reduzieren und in Folge dessen so genannte Filterblasen (Filter Bubbles) entstehen \textcolor{red}{(Haim, Graefe \& Brosius, 2017)}.  
Visualisierung
TreeMaps sind eine Visualisierungsform, bei der verschachtelte Rechtecke verwendet werden, um große Mengen monohierarchischer, baumstrukturierter Daten mit hoher Effizienz und hoher Lesbarkeit anzuzeigen \textcolor{red}{(Gasteier, Krois \& Hrdina, 2013)}.
Sie bieten den entscheidenden Vorteil, dass große Teile des gesamten Suchraums, also alle verfügbaren Alternativen, strukturiert und übersichtlich abgebildet werden können. Dadurch können die relevanten Optionen besser in ihren Kontext eingeordnet werden. Dies trägt zu einem besseren Verständnis über die verfügbaren Optionen bei und erleichtert die Navigation im Suchraum \textcolor{red}{(Tintarev \& Masthoff, 2007)}. Durch den EInsatz verschiedener Farben und Deckkraftswerte können zusätzliche Attribute dargestellt werden oder eine Akzentuierung vorgenommen werden. Weitere Gestaltungsmöglichkeiten ergeben sich durch die Verwendugn verschiedener Schriftarten, Schriftgrößen oder Rahmenarten \textcolor{red}{(Richthammer \& Pernul, 2016)}.
Die Visualisierung des Suchraumes in dem sich der Benutzer bewegt trägt dazu bei, dass der Benutzer die Gründe hinter den Empfehlungen besser nachvollziehen kann \textcolor{red}{(Smyth nach Richthammer \& Pernul, 2016)}. Außerdem fungiert sie als Ausgangspunkt für gezielte Interaktion und Manipulation durch den Benutzer. 


\subsection{Beschreibung des Problems}
\textcolor{red}{Beschreibung des konkreten Problems, das, durch die oben identifizierte Lücke in der wissenschaftlichen Literatur motiviert, in dieser Arbeit angegangen werden soll.}

---------------------------------------------------------

\begin{itemize}
\item Visualisierung des gesamten Informationsraumes mit TreeMap Visualisierung --> Transparenz, Diversität (Anti Filter Bubble), Kontext
\end{itemize}
\begin{itemize}
\item Anwenden auf neue Domäne: Online News, neue Visualisierung (TreeMap) die sich gut eignet um große Datenmengen zu visualisieren
\end{itemize}


---------------------------------------------------------

\subsection{Ziele}
\textcolor{red}{Identifikation der zentralen Forschungsfrage (oder Forschungsfragen), welche sich aus dem zuvor beschriebenen Problem ableitet. Dabei Erläuterung des Ziels der geplanten Arbeit, d.h. wie die Forschungsfrage im Rahmen dieser Arbeit konkret behandelt werden soll.}

---------------------------------------------------------
\begin{itemize}
\item Visualisierung des Gesamten Informationsraumes in 2D mit TreeMap
\item Transparenz, Übersicht, Interaktion
\item Diversität (Filter Bubble, Echo Chamber)
\end{itemize}

\begin{enumerate}
\item Sind TreeMaps geeignet um Online News mit Empfehlungen übersichtlich zu visualisieren?
\item Trägt die Visualisierung & Interaktion damit dazu bei Transparenz der Empfehlungen zu schaffen?
\item Diversere Nachrichten durch Visualisierung des gesamten Informationsraumes? 
\end{enumerate}

---------------------------------------------------------
\section{Stand der Forschung}
\textcolor{red}{Vertiefung der oben kurz angeschnittenen Hintergründe. Dabei Beschreibung der Ansätze, die es bereits für die Lösung des Problems gibt. Besonders sollten Publikationen berücksichtigt werden, die zu Ihrer Lösung am meisten beitragen können und die Sie wahrscheinlich stark bei Ihrer Arbeit beeinflussen werden. Neben einer guten thematischen Passung, sollten zitierte Quellen möglichst aktuell und bestenfalls bereits sehr häufig zitiert worden sein.}

\textcolor{red}{Sie sollten an dieser Stelle mindestens fünf einflussreiche Publikationen zitieren. Zitationen können entweder direkt im Text (z.B. \enquote{\citet{Zobel:2004} zeigten \emph{XY}}) oder als nachgestellte Quellen (z.B. \enquote{\emph{XY} wurde bereits gezeigt \citep{Zobel:2004}.}) erscheinen. Alle zitierten Publikationen müssen korrekt und vollständig im Literaturverzeichnis aufgeführt werden. Für weitere Hinweise und formale Richtlinien, siehe auch den vom Lehrstuhl zur Verfügung gestellten \emph{Primer}\footnote{\url{http://interactivesystems.info/theses/howto}}.}

\textcolor{red}{Am Ende dieses Abschnittes sollten Sie herausarbeiten, weshalb die bisherigen Ansätze nicht ausreichend sind: Welche Defizite weisen diese Arbeiten auf? Was wurde nicht untersucht? Worin unterscheidet sich Ihre geplante Arbeit von den bisherigen?}

---------------------------------------------------------

\begin{itemize}
\item  Kunkel, J., Loepp, B., Ziegler, J. (2015) - 3D Visualisierung zur Eingabe von Präferenzen in Empfehlungssystemen
\item He, C., Parra, D., Verbert, K. (2016) - Interactive Recommender Systems - A Survey of the State of the Art and Future Research Challenges and Opportunities (Übersicht über bisherige Ansätze)
\item Richthammer (2016): Interactive Erkundung des Empfehlungsraums, Visualisierung mit TreeMaps, Bewertung von Filmen zu Beginn, Collaborative Filtering --> 
\end{itemize}


Nicht untersucht: Treemaps, Online News

---------------------------------------------------------

\section{Konzept}
\textcolor{red}{Grobe Skizzierung des Konzepts, welches zur Lösung des Problems (bzw. zum Erreichen der oben genannten Ziele) angedacht ist. Die entscheidenden Punkte des Konzepts sollten dabei schon konkreter benannt werden. Allerdings können Details (insbesondere bezüglich der Implementierung) an dieser Stelle noch vernachlässigt werden. Wichtig ist, dass ein klares Bild von der umzusetzenden Lösung entsteht und daraus ersichtlich wird, wie das zuvor genannte Problem effektiv angegangen werden kann.}

---------------------------------------------------------

\begin{itemize}
\item Implementierung einer prototypischen Webanwendung (JSF-Framework, JS)
\item Visualisierung des gesamten Informationsraums und der Empfehlungen als TreeMap Interaktions- / Manipulationsmöglichkeiten (Zoom, Filter, Rotation, Bewegung, Items add/remove), …)
\item Ansicht und Manipulation des Nutzerprofils
\end{itemize}

Movie Landkarte

\begin{enumerate}
\item Ähnlichkeit zwischen Items berechnen (auf Basis der Nutzerbewertungen)  Matrix-Factorization
\item Zweidimensionale, nutzerübergreifende Anordnung der Items --> MDS
\item Clustering (populärstes Item als Vertreter abgebildet)
\end{enumerate}

--> Frage: Soll bei mir das gleiche Vorgehen genutzt werden?

---------------------------------------------------------

\section{Realisierung}
\textcolor{red}{Kurze und eher allgemeine Vorstellung, wie das Konzept (insbesondere in technischer Hinsicht) umgesetzt werden soll.}

---------------------------------------------------------
\begin{itemize}
\item JSF Framework (Java) --> Grundgerüst steht (Hello World)
\item Datenbank mit News Artikeln
\item D3.js (plotly.js)?
\item Interaktionselemente realisieren
\end{itemize}


Beispiele:

\begin{itemize}
\item \href{https://bost.ocks.org/mike/treemap/}{https://bost.ocks.org/mike/treemap/}
\item \href{http://bl.ocks.org/ganeshv/6a8e9ada3ab7f2d88022}{http://bl.ocks.org/ganeshv/6a8e9ada3ab7f2d88022}
\item \href{https://d3indepth.com/layouts/}{(https://d3indepth.com/layouts/)}
\end{itemize}


---------------------------------------------------------

\section{Validierung}
\textcolor{red}{Beschreibung, wie geplant ist nachzuweisen, dass die Lösung tut, was sie soll. Es gilt auch zu erläutern, welche Methoden angedacht sind und wie die Lösung validiert bzw. evaluiert werden soll. Dabei ist es immer gut, auf etablierte Verfahren zurückzugreifen (und andernfalls hier zu begründen, warum man es nicht tut). Bei einer empirischen Untersuchung sollte die Überschrift besser \enquote{Evaluation} lauten. Je nach Art und Ausrichtung der Arbeit kann dieser Abschnitt ggf. entfallen.}

---------------------------------------------------------

Studie entsprechend (\textcolor{red}{Kunkel, Loepp, Ziegler, 2015})  Anwendung des Ansatzes auf neue Domäne & andere Visualisierung aber ähnliche Validierung sinnvoll

Laborstudie -->  Anwendung im Webbrowser nutzen

\textbf{Fragebögen}

\begin{itemize}
\item Demografie
\item Domänenkenntnis
\end{itemize}

--> allg. Fragebögen

\begin{itemize}
\item Empfehlungsqualität
\item Effizienz des Systems
\item Transparenz der Empfehlungsgenerierung
\end{itemize}


--> Items von (\textcolor{red}{ Knijnenburg et al. (2012) und Pu et al. (2011)})

\begin{itemize}
\item Möglichkeiten Kontrolle auszuüben
\item Nutzungsabsicht
\item Systemspezifische Funktionalitäten (Verständlichkeit der Landschaft, Nützlichkeit des Modellierungswerkzeugs, …)
\end{itemize}


--> selbstgenerierte Items (\textcolor{red}{Kunkel, Loepp, Ziegler, 2015})

\begin{itemize}
\item Usability
\item User Experience
\end{itemize}


-->  System Usability Scale (\textcolor{red}{(SUS; Brooke 1996)}) und den User Experience Questionnaire (\textcolor{red}{(UEQ; Laugwitz et al. 2008)}) 

\begin{itemize}
\item Genereller Nutzen des Systems
\item Motivation & Spaß bei der Verwendung
\item Interesse der Nutzer
\end{itemize}


--> Intrinsic Motivation Inventory Scale (\textcolor{red}{(IMI; Ryan 1982)}). Bis auf bei UEQ (7-stufige bipolare Skala) und IMI (7-stufige Likert-Skala) wurden sämtliche Items auf einer positiven 5-stufigen Likert-Skala (1−5) erhoben

\begin{itemize}
\item Interaktionsverhalten
\item Aufgabendauer
\end{itemize}


--> Logdateien

---------------------------------------------------------


\bibliographystyle{apacite}

%\bibliography{Literatur}

\end{document}